\section{Lux Network Integration}

\label{sec:lux}
% ============================================================================

Base integrates with the Lux Network at three chain layers, each providing a
distinct sovereignty primitive~\cite{luxwhitepaper, luxtrust, luxapp}.

\subsection{K-Chain: Key Policy}

The Lux K-Chain (Key Chain) manages cryptographic key policies. Base registers
its master key with the K-Chain, enabling:

\begin{itemize}[leftmargin=1.1em]
    \item \textbf{Key rotation}: The K-Chain enforces rotation schedules. Base
    re-encrypts field-level DEKs when the KEK rotates, without re-encrypting
    the underlying data (envelope encryption).
    \item \textbf{Key escrow}: A user can designate a recovery agent on the
    K-Chain. If the user loses their master key, the recovery agent can
    authorize re-derivation via threshold signature.
    \item \textbf{Key revocation}: Compromised keys are revoked on-chain.
    All Base instances observing the K-Chain stop accepting operations signed
    by the revoked key within one block (sub-second).
\end{itemize}

\subsection{I-Chain: Identity}

The Lux I-Chain (Identity Chain) provides decentralized identifiers (DIDs)
and verifiable credentials. Base uses I-Chain identities for:

\begin{itemize}[leftmargin=1.1em]
    \item \textbf{User authentication}: Base accepts I-Chain DIDs as a first-class
    auth method alongside Hanzo IAM OIDC tokens. A user with an I-Chain DID
    can authenticate to any Base instance without a centralized identity
    provider.
    \item \textbf{Collection ownership}: Collections can be bound to an I-Chain
    DID rather than a Hanzo org. This enables personal data stores that are
    not scoped to any organization.
    \item \textbf{Credential verification}: Access rules can reference verifiable
    credentials issued on I-Chain (e.g., ``only users with a verified
    healthcare credential can access this collection'').
\end{itemize}

\subsection{T-Chain: Threshold Controls}

The Lux T-Chain (Threshold Chain) implements threshold cryptography for
multi-party authorization. Base uses T-Chain for:

\begin{itemize}[leftmargin=1.1em]
    \item \textbf{Threshold write attestation}: For collections requiring
    multi-party writes (e.g., financial records, legal documents), T-Chain
    enforces $k$-of-$n$ signature requirements before the write is committed.
    \item \textbf{Threshold decryption}: Encrypted fields can be configured to
    require $k$-of-$n$ T-Chain participants to authorize decryption, preventing
    any single party from accessing the plaintext.
    \item \textbf{Threshold key recovery}: Master key recovery requires
    $k$-of-$n$ designated recovery agents to cooperate, preventing both
    single-point-of-failure and unilateral recovery.
\end{itemize}

% ============================================================================
